\documentclass[11pt]{article}

\usepackage[T1]{fontenc}
\usepackage{amsmath,amssymb,amsthm}
\usepackage[margin=1in]{geometry}
\usepackage{xcolor}
\usepackage{hyperref}
\hypersetup{colorlinks=true,linkcolor=blue,urlcolor=blue,citecolor=blue}

\theoremstyle{plain}
\newtheorem{theorem}{Theorem}[section]
\newtheorem{lemma}[theorem]{Lemma}
\newtheorem{proposition}[theorem]{Proposition}
\newtheorem{corollary}[theorem]{Corollary}
\newtheorem{conjecture}[theorem]{Conjecture}
\theoremstyle{definition}
\newtheorem{definition}[theorem]{Definition}
\theoremstyle{remark}
\newtheorem{remark}[theorem]{Remark}

\newcommand{\R}{\mathbb R}
\newcommand{\C}{\mathbb C}
\newcommand{\Z}{\mathbb Z}
\newcommand{\CAZ}{\mathrm{CAZ}}
\newcommand{\dist}{\operatorname{dist}}
\newcommand{\re}{\operatorname{Re}}
\newcommand{\iC}{\iota_{C_\tau}}
\newcommand{\El}{\mathcal E_n}

\title{Convergence of iterated SOCP refinement for CAZAC feasibility}
\author{Jeffery Kline}
\date{July 2026}

\begin{document}
\maketitle

\begin{abstract}
Constant-amplitude zero-autocorrelation (CAZAC) sequences satisfy a
constant-modulus condition in time and a flat-spectrum condition in
frequency simultaneously, a nonconvex intersection of two tori. Replacing
the equality constraints with convex inequalities (bounded amplitude,
bounded spectral energy) yields a compact convex semialgebraic body
$C_\tau$ on which CAZAC systems are exactly the maximal-norm points; the
natural algorithm is to iterate a second-order cone program (SOCP) whose
objective direction resets to the previous optimizer. We give a complete
convergence theory for this iteration. Theorem A (unconditional): every
measurable selection from the iteration has finite length and converges to
a single fixed point, via the Attouch--Bolte--Svaiter Kurdyka--{\L}ojasiewicz
(KL) descent framework, applied here to a set-valued argmax update rather
than a proximal step. The same argument, since it never uses CAZAC-specific
structure beyond compactness, convexity, and semialgebraicity, transfers
verbatim to a second example, the elliptope $\El=\{X\succeq0: X_{ii}=1\}$ ---
Felzenszwalb--Klivans--Paul's (2021) own worked case, where their
finiteness-dependent single-point-convergence mechanism does not apply and
their own example stalls at set-level convergence; we prove single-point
convergence there unconditionally and verify numerically (four trajectories
on $\mathcal E_5$) that generic trajectories escape the fixed-point continuum
onto a rank-1 point, a phenomenon we leave as Conjecture (Q7). Theorem B
(unconditional, obstruction): non-CAZAC fixed points exist and some are
exposed (unique maximizers of their own functional), so no deterministic
selection rule can guarantee convergence to the CAZAC set; the correct
target class of theorem is probabilistic over initializations. Theorem C
(proved conditional on a checkable regularity hypothesis (R), unconditional
at Zadoff--Chu points of squarefree length): near a regular CAZAC point the
$\ell^1$ constraint slack has \emph{sharp linear}, not merely quadratic,
growth; the KL exponent is exactly $0$; and the iteration lands on the
point \emph{exactly} after finitely many steps. We give the complete proof
of part (c) (finite exact arrival), and prove hypothesis (R) at
Zadoff--Chu points for all $n$: $\dim K=\max\{d: d^2\mid n\}$ regardless
of parity, so (R) holds iff $n$ is squarefree. Beyond Zadoff--Chu points we
report a partial result at Bj\"orck's quadratic-phase construction at prime
$n$: an affine-in-Legendre-symbol closure argument proves the dependency
space meets the decimation-symmetry-invariant subspace only in the constant
direction, reducing regularity there to one unproved, numerically-confirmed
(44 primes) determinant-nonvanishing lemma. We close with a precise
statement of the remaining open problems (Q1--Q7) and a tiered
verification ledger disclosing exactly which claims are proved, proved
conditionally, numerically verified, or conjectural.
\end{abstract}

\section{Setup}

Fix $n\ge1$ and $N\ge1$. Let $F$ denote the unnormalized discrete Fourier
transform on $\C^n$, $F(u)(k)=\sum_m u(m)e^{-2\pi ikm/n}$, so Parseval's
identity reads $\|Fu\|^2=n\|u\|^2$. The state space is
$x=(x_0,\dots,x_{N-1})\in(\C^n)^N\cong\R^{2Nn}$, with real inner product
$\langle u,v\rangle_\R=\re\langle u,v\rangle$ and induced norm
$\|x\|^2=\sum_{j,k}|x_j(k)|^2$. Write the \emph{coupled spectrum}
$S_x(k)=\sum_{j=0}^{N-1}|F(x_j)(k)|^2$.

\begin{definition}[Coupled CAZAC system]
$x$ is a coupled CAZAC system if $|x_j(k)|=1$ for all $j,k$ and
$S_x(k)=Nn$ for all $k$. Write $\CAZ(n,N)$ for this set; for $N=1$ this is
the classical CAZAC condition. $\CAZ(n,N)$ is nonempty for every $n\ge1$,
$N\ge1$: a Zadoff--Chu sequence exists for every length $n$ \cite{Chu72},
and $N$ copies of one satisfy the coupled constraint.
\end{definition}

For $\tau\ge1$, the convex relaxation is
$$C_\tau=\bigl\{x\in(\C^n)^N: |x_j(k)|\le\tau\ \forall j,k,\ \ S_x(k)\le Nn\ \forall k\bigr\}.$$
$C_\tau$ is compact (amplitude bounds), convex, semialgebraic, and
second-order-cone representable (each constraint is a Euclidean-norm bound
on a linear image of $x$); it is invariant under global phase rotation of
each $x_j$, cyclic shifts, and (for the coupled constraint) any
transformation preserving $S_x$.

\textbf{Iteration.} Given $x^0\ne0$, the algorithm is \emph{iterated linear
optimization}:
\begin{equation}
x^{\ell+1}\in\operatorname*{arg\,max}_{x\in C_\tau}\langle x^\ell,x\rangle_\R. \tag{$\star$}
\end{equation}
The argmax may be non-unique; $(\star)$ permits any measurable selection.
The paper's practical schedule presolves at $\tau=2+1/\ell$ for
$\ell=1,\dots,10$, then fixes $\tau=1$ for the main phase; all theorems
below concern the main phase at fixed $\tau\ge1$ (in particular $\tau=1$),
and hold for any measurable selection rule, so the finitely many presolve
steps do not affect any conclusion.

\begin{definition}[Fixed point]
$x^*$ is a fixed point of $(\star)$ if
$x^*\in\arg\max_{x\in C_\tau}\langle x^*,x\rangle_\R$, equivalently
$\sigma_{C_\tau}(x^*)=\|x^*\|^2$, where
$\sigma_C(r)=\max_{x\in C}\langle r,x\rangle_\R$ is the support function of $C$.
\end{definition}

Three increasingly strong convergence questions organize everything that
follows: (1) does $x^\ell$ converge to a single point, not merely have
limit points; (2) is that limit CAZAC; (3) at what rate.

\section{The iteration as a generalized power method}

\begin{proposition}\label{prop:1}
With $f(x)=\tfrac12\|x\|^2$ (convex, $\nabla f(x)=x$), $(\star)$ is exactly
$x^{\ell+1}\in\arg\max_{x\in C_\tau}\langle\nabla f(x^\ell),x\rangle_\R$: the
conditional-gradient / generalized-power-method update for maximizing the
convex function $f$ over the compact convex set $C_\tau$ \cite{JNRS10},
also called iterated linear optimization \cite{FKP21}.
\end{proposition}
\begin{proof}
Immediate from the definitions.
\end{proof}

\begin{proposition}[Max-norm points are exactly CAZAC points]\label{prop:2}
For every $x\in C_\tau$ ($\tau\ge1$), $\|x\|^2\le Nn$. For $\tau=1$,
equality holds iff $x\in\CAZ(n,N)$. The maximum is attained for every
$n,N,\tau$.
\end{proposition}
\begin{proof}
By Parseval, $\|x\|^2=\sum_j\|x_j\|^2=\tfrac1n\sum_k S_x(k)\le\tfrac1n\cdot n\cdot Nn=Nn$,
with equality iff $S_x(k)=Nn$ for every $k$. Separately, at $\tau=1$,
$\|x\|^2=\sum_{j,k}|x_j(k)|^2\le Nn$ with equality iff every $|x_j(k)|=1$.
Both chains tight simultaneously forces both CAZAC conditions; conversely a
CAZAC system attains $Nn$. Attainment: $N$ copies of a Zadoff--Chu sequence
of length $n$ lie in $C_1\subseteq C_\tau$ with norm$^2=Nn$.
\end{proof}

This reframes the nonconvex feasibility problem ``find a CAZAC system'' as
the convex maximization $\max_{x\in C_1}\|x\|^2$, whose optimal value $Nn$
is known a priori: an iterate is $\delta$-suboptimal iff $Nn-\|x^\ell\|^2=\delta$.

\begin{proposition}[Monotone ascent, square-summable increments]\label{prop:3}
Along the iteration at fixed $\tau$, with $x^\ell\ne0$:
\begin{enumerate}
\item[\rm(1)] $\langle x^\ell,x^{\ell+1}\rangle_\R\ge\|x^\ell\|^2$;
\item[\rm(2)] $\|x^{\ell+1}\|\ge\|x^\ell\|$, and $\|x^\ell\|\uparrow\rho\le\sqrt{Nn}$;
\item[\rm(3)] $\|x^{\ell+1}-x^\ell\|^2\le\|x^{\ell+1}\|^2-\|x^\ell\|^2$, hence
$\sum_{\ell\ge0}\|x^{\ell+1}-x^\ell\|^2\le Nn-\|x^0\|^2$ and
$\|x^{\ell+1}-x^\ell\|\to0$.
\end{enumerate}
\end{proposition}
\begin{proof}
(1) $x^\ell\in C_\tau$ is feasible for the SOCP whose optimizer is
$x^{\ell+1}$, so the optimum is at least $\langle x^\ell,x^\ell\rangle_\R=\|x^\ell\|^2$.
(2) Cauchy--Schwarz: $\|x^\ell\|\|x^{\ell+1}\|\ge\langle x^\ell,x^{\ell+1}\rangle_\R\ge\|x^\ell\|^2$;
divide by $\|x^\ell\|$; monotone and bounded above by Prop.~\ref{prop:2}.
(3) $\|x^{\ell+1}-x^\ell\|^2=\|x^{\ell+1}\|^2+\|x^\ell\|^2-2\langle x^\ell,x^{\ell+1}\rangle_\R\le\|x^{\ell+1}\|^2-\|x^\ell\|^2$
by (1); telescope and apply (2).
\end{proof}

Square-summable increments alone do not imply convergence of the full
sequence (harmonic-type sums can wander); closing exactly this gap is
Theorem A.

\begin{proposition}[Limit points are fixed points; CAZAC points are fixed points]\label{prop:4}
Every limit point of $(x^\ell)$ is a fixed point of $(\star)$. Every
$x^*\in\CAZ(n,N)$ is a fixed point of the $\tau=1$ iteration, and moreover
a global maximizer of its own functional:
$\arg\max_{x\in C_1}\langle x^*,x\rangle_\R\ni x^*$ with optimal value $Nn$.
\end{proposition}
\begin{proof}
Let $x^{\ell_j}\to\bar x$. By Prop.~\ref{prop:3}(3), $x^{\ell_j+1}\to\bar x$
as well. $\sigma_{C_\tau}$ is continuous (finite convex on $\R^{2Nn}$), and
$\langle x^{\ell_j},x^{\ell_j+1}\rangle_\R=\sigma_{C_\tau}(x^{\ell_j})$ by
definition of the update; pass to the limit. For the second claim, any
$y\in C_1$ satisfies $\langle x^*,y\rangle_\R\le\|x^*\|\|y\|\le\sqrt{Nn}\sqrt{Nn}=Nn=\|x^*\|^2$
by Cauchy--Schwarz and Prop.~\ref{prop:2}, with equality at $y=x^*$.
\end{proof}

The same Cauchy--Schwarz argument shows any maximal-norm point of \emph{any}
compact convex set is a fixed point of iterated linear optimization;
nothing CAZAC-specific is used beyond Prop.~\ref{prop:2} itself --- a fact
that drives the elliptope corollary in Section~\ref{sec:elliptope}.

\section{Theorem A --- global convergence}

\begin{theorem}[Theorem A --- Proved, unconditional]\label{thm:A}
Fix $\tau\ge1$. For any $x^0\ne0$ and any measurable selection
$x^{\ell+1}\in\arg\max_{x\in C_\tau}\langle x^\ell,x\rangle_\R$, the sequence
$(x^\ell)$ has finite length, $\sum_\ell\|x^{\ell+1}-x^\ell\|<\infty$, and
converges to a single fixed point $\bar x\in C_\tau$.
\end{theorem}
\begin{proof}
Apply the abstract descent theorem of Attouch, Bolte, and Svaiter
\cite[Thm.~2.9]{ABS13} to $F(x)=-\tfrac12\|x\|^2+\iC$ (so minimizing $F$ is
maximizing $\|x\|^2/2$ on $C_\tau$):
\begin{itemize}
\item \textbf{Sufficient decrease} ($a=\tfrac12$): exactly Prop.~\ref{prop:3}(3).
\item \textbf{Relative error} ($b=1$), the structural surprise: first-order
optimality of the argmax step gives $x^\ell\in N_{C_\tau}(x^{\ell+1})$, the
normal cone at the new iterate. Since $\partial F(x^{\ell+1})=-x^{\ell+1}+N_{C_\tau}(x^{\ell+1})$
(sum rule for $C^1$ plus indicator of a convex set), the vector
$w=x^\ell-x^{\ell+1}$ lies in $\partial F(x^{\ell+1})$, and
$\|w\|=\|x^{\ell+1}-x^\ell\|$ \emph{exactly}: the subgradient residual is
literally the step. This identity is specific to $f=\tfrac12\|x\|^2$; for a
generic convex objective the relative-error condition would be a genuine
additional hypothesis, but here it is free.
\item \textbf{Compactness/continuity}: iterates lie in compact $C_\tau$;
$F$ is continuous there.
\item \textbf{KL property}: $C_\tau$ is defined by polynomial (in)equalities
and $F$ is polynomial plus the indicator of a semialgebraic set, hence
semialgebraic, hence KL \cite{BDLS07}.
\end{itemize}
Critical points of $F$ are exactly the fixed points of $(\star)$
(Sec.~2). The theorem then gives finite length and convergence to a single
critical point.
\end{proof}

The theorem is silent on \emph{which} fixed point the limit is; that is
exactly where Theorem B enters. Selection-independence, and the fact that
the presolve is finitely many extra solves, mean the conclusion applies
verbatim to any concrete implementation of $(\star)$.

\subsection{A second worked example: the elliptope}\label{sec:elliptope}

The proof of Theorem~\ref{thm:A} uses nothing about $C_\tau$ beyond
compactness and convexity, except at the KL step, which needs
semialgebraicity. Stated in general: for \emph{any} compact convex
$C\subset\R^d$ and the iteration
$x^{\ell+1}\in\arg\max_{x\in C}\langle x^\ell,x\rangle$, sufficient decrease
(Prop.~\ref{prop:3}(3)'s proof uses only feasibility of $x^\ell$ for the
next problem and Cauchy--Schwarz) and the relative-error identity (the
normal-cone argument above, using only $x^\ell\in N_C(x^{\ell+1})$ and
$\nabla(\tfrac12\|x\|^2)=x$) hold unconditionally; compactness gives
continuity. If $C$ is also semialgebraic, KL holds \cite{BDLS07} and
Theorem A's conclusion --- convergence to a single fixed point, not merely a
connected limit set --- follows for that $C$ with no CAZAC-specific input.

We exercise this on a genuinely different $C$: the \textbf{elliptope}
$$\El=\{X\in\R^{n\times n}: X=X^\top,\ X\succeq0,\ X_{ii}=1\ \forall i\},$$
the set of $n\times n$ correlation matrices, with iteration
$X^{\ell+1}\in\arg\max_{X\in\El}\langle X^\ell,X\rangle$
($\langle A,B\rangle=\operatorname{tr}(AB)$). This is Felzenszwalb, Klivans,
and Paul's own example \cite{FKP21}: the one case in their paper where they
concede the fixed-point set is a continuum for $n>3$, and where their own
worked instance stops at set-level convergence.

\begin{proposition}[Proved]\label{prop:ell1}
$\El$ is compact, convex, and semialgebraic.
\end{proposition}
\begin{proof}
Convexity of the constraints is immediate. Compactness: $X\succeq0$ with
$X_{ii}=1$ gives $|X_{ij}|\le\sqrt{X_{ii}X_{jj}}=1$ by Cauchy--Schwarz on
the associated Gram vectors, so $\El$ is bounded, and closed as an
intersection of closed sets, hence compact. Semialgebraicity: for
symmetric $X$, $X\succeq0$ is equivalent by Sylvester's criterion to the
finite conjunction of polynomial inequalities ``every principal minor of
$X$ is $\ge0$''; the diagonal constraints are polynomial equalities. So
$\El$ is a basic semialgebraic set (a spectrahedron).
\end{proof}

\begin{proposition}[Proved]\label{prop:ell2}
Sufficient decrease, relative error, and compactness/continuity transfer
verbatim from Prop.~\ref{prop:3}(3) and the normal-cone argument of
Theorem~\ref{thm:A}, replacing ``$C_\tau$'' by ``$\El$'' and
$\langle\cdot,\cdot\rangle_\R$ by $\operatorname{tr}(\cdot\,\cdot)$.
\end{proposition}
\begin{proof}
Neither source proof uses any CAZAC structure; both use only convexity,
Cauchy--Schwarz, and the normal-cone characterization of the argmax, all of
which hold verbatim on $\El\subset\R^{\binom{n+1}2}$.
\end{proof}

\begin{proposition}[Proved]\label{prop:ell3}
The fixed-point set of the elliptope iteration is a genuine continuum, not
a singleton: $X^*=I_n$ is a fixed point whose optimal face at that
direction is all of $\El$.
\end{proposition}
\begin{proof}
For every $X\in\El$, $\operatorname{tr}(I_nX)=\sum_iX_{ii}=n$ identically ---
the linear functional $\langle I_n,\cdot\rangle$ is constant on all of
$\El$. Hence $\arg\max_{X\in\El}\langle I_n,X\rangle=\El$ itself: every
feasible point, including every rank-1 correlation matrix $vv^\top$
($v_i=\pm1$), is simultaneously optimal. This instantiates, with a
machine-exact identity, the general fact already noted by Felzenszwalb,
Klivans, and Paul \cite{FKP21}, who state that $\mathcal L_n$ (their name
for $\El$) has infinitely many fixed points for $n>3$.
\end{proof}

\begin{proposition}[Verified numerically, $n=5$]\label{prop:ell4}
Running the iteration on $\mathcal E_5$ from four starting points ---
$X^0=I_5$; $I_5$ plus a small random perturbation renormalized to unit
diagonal; and two independent random correlation matrices --- using an
interior-point SDP solver (CLARABEL via cvxpy), all four trajectories show
$\|X^{\ell+1}-X^\ell\|_F\to0$ (to $\sim10^{-15}$ by iteration $15$) and
settle on a single matrix. Three of the four escape the fixed-point
continuum entirely and converge to a rank-1 correlation matrix $vv^\top$
(eigenvalues $(0,0,0,0,5)$ to machine precision); the trajectory started at
$X^0=I_5$ is an edge case in which the solver's own deterministic tie-break
returns $I_5$ again, a valid fixed point, not a counterexample. Script:
{\tt experiments/verify\_elliptope\_convergence.py} in the companion
repository \cite{CazacAlgorithmRepo}.
\end{proposition}

\begin{theorem}[Elliptope convergence --- tiers as stated]\label{thm:ell}
Because $\El$ is compact, convex, and semialgebraic (Prop.~\ref{prop:ell1}),
Theorem A's general form applies to it directly: the iterated-linear-
optimization sequence on the elliptope converges to a single fixed point
for every initialization, \textbf{despite the fixed-point set being an
infinite continuum} (Prop.~\ref{prop:ell3}) --- \textbf{Proved}. That
generic nearby trajectories escape the continuum and settle onto a single
point rather than merely having a connected limit set is
\textbf{verified numerically} (Prop.~\ref{prop:ell4}, $n=5$, four starts),
not established for every $n$ or every trajectory. Whether \emph{every}
trajectory converges specifically to a rank-1 extreme point is
\textbf{open} (Conjecture, filed as Q7 below).
\end{theorem}

\begin{remark}[Relation to Felzenszwalb--Klivans--Paul]
FKP's Theorem~1 (limit points exist, are fixed points, form a connected
set) is the same starting point as Prop.~\ref{prop:4} here, but their only
route to single-point convergence is an unnumbered remark that a connected
\emph{finite} fixed-point set is a singleton --- structurally inapplicable
whenever the fixed-point set is a positive-dimensional continuum, which is
exactly the situation both here (the $N=1$ CAZAC phase circle) and in their
own \S4.2 example. A full read of the paper's body, not merely its
abstract, confirms no alternate route exists in their text: the proof of
their Theorem~1 only ever derives connectedness and never separately
invokes finiteness. Theorem~\ref{thm:ell} closes exactly this gap via the
KL/finite-length route \cite{ABS13} rather than the connectedness-plus-
finiteness route, and is therefore not subsumed by \cite{FKP21} even in the
positive-dimensional-face case.
\end{remark}

\section{Theorem B --- obstruction: the limit need not be CAZAC}

\begin{theorem}[Theorem B --- Proved, unconditional]\label{thm:B}
\begin{enumerate}
\item[\rm(a)] The fixed points of $(\star)$ at $\tau=1$ with every
amplitude constraint slack are exactly the points whose Fourier transform
has modulus $\sqrt n$ on some proper subset $K\subsetneq\{0,\dots,n-1\}$ of
frequencies and vanishes off $K$, with time-domain modulus strictly less
than $1$ everywhere; such a point has $\|x\|^2=|K|$.
\item[\rm(b)] There exist non-CAZAC fixed points that are the
\textbf{unique} maximizer of their own linear functional over $C_1$ ---
exposed points. For $n=2$: $x^*=(1,t)$, $t=\sqrt2-1$, is such a point; for
$N=2$, its diagonal doubling $((1,t),(1,t))$ is such a point.
\item[\rm(c)] Consequently no deterministic convergence-to-CAZAC theorem
is possible: every measurable selection rule in $(\star)$ retains the
point of (b) exactly, once reached.
\end{enumerate}
\end{theorem}
\begin{proof}[Proof of (a), $N=1$]
By Prop.~\ref{prop:1}, $x^*\ne0$ is a fixed point iff it is a KKT point of
$\max_y\langle x^*,y\rangle_\R$ subject to $g_m(y)=|y(m)|^2-1\le0$ and
$h_k(y)=|\widehat y(k)|^2-n\le0$; this is a convex program (linear
objective, convex constraints) and Slater's condition holds at $y=0$
($g_m(0)=-1<0$, $h_k(0)=-n<0$ strictly), so KKT is necessary \emph{and}
sufficient for optimality of $x^*$. The real gradients are
$\nabla g_m(y)=2y(m)e_m$ and $\nabla h_k(y)=2F^*E_kFy$ ($E_k=e_ke_k^*$), so
KKT at $x^*$ reads $x^*=\Lambda x^*+F^*MFx^*$ for diagonal $\Lambda,M\ge0$,
complementary to the active sets. If every amplitude constraint is slack
then $\Lambda=0$, so $x^*=F^*MFx^*$; applying $F$ and using $FF^*=nI$
gives $\widehat x^*=nM\widehat x^*$ coordinatewise, i.e.\ for each $k$
either $\widehat x^*(k)=0$ or $\mu_k=1/n$. Let
$K=\{k:\widehat x^*(k)\ne0\}\ne\varnothing$; for $k\in K$ complementary
slackness with $\mu_k=1/n>0$ forces $|\widehat x^*(k)|^2=n$. Parseval gives
$\|x^*\|^2=|K|$; if $|K|=n$ this forces $\|x^*\|^2=n$, hence (Prop.~2)
$x^*\in\CAZ(n,1)$, contradicting amplitude-slackness. So $K\subsetneq\{0,\dots,n-1\}$.
Conversely, for such $x^*$ and any $y\in C_1$,
$\langle x^*,y\rangle_\R=\tfrac1n\sum_{k\in K}\re(\widehat x^*(k)\overline{\widehat y(k)})\le\tfrac1n\sum_{k\in K}\sqrt n\sqrt n=|K|=\|x^*\|^2$,
attained at $y=x^*$ --- a direct support-function argument, no KKT needed
for sufficiency.
\end{proof}
\begin{proof}[Proof of (b), $n=2,N=1$]
For $y=(y_0,y_1)\in\C^2$ feasible in $C_1$ (so $|y_m|\le1$ and
$|F(y)(0)|^2+|F(y)(1)|^2\le2$, hence $|F(y)(0)|\le\sqrt2$),
$$\langle x^*,y\rangle_\R=(1-t)\,\re y(0)+t\,\re F(y)(0)\le(1-t)+t\sqrt2=1+t^2=\|x^*\|^2,$$
with equality iff $y(0)=1$ and $F(y)(0)=\sqrt2$, which forces
$y(1)=\sqrt2-1=t$, i.e.\ $y=x^*$ --- a genuine separating functional
exposing $x^*$ as a vertex-like point of $C_1$ even though
$x^*\notin\CAZ(2,1)$ ($\|x^*\|^2=4-2\sqrt2<2$). The $N=2$ doubling repeats
the argument coordinatewise. Part (c) follows immediately from the
existence of one such retained point.
\end{proof}

\begin{remark}
Theorem A guarantees convergence to \emph{some} fixed point; Theorem B
shows the fixed-point set genuinely contains non-CAZAC points that trap
every selection rule. The remaining content of ``the algorithm works'' is a
probabilistic/basin-of-attraction statement (Q6 below), not a deterministic
one.
\end{remark}

\section{Theorem C --- finite exact termination at regular CAZAC limits}

Call $x^*\in\CAZ(n,N)$ \textbf{regular} (hypothesis (R)) if the only
linear dependency among the gradients of the active constraints at $x^*$
is the one forced by Parseval's identity; equivalently the dependency
space $K$ (Section~\ref{sec:zc} makes this precise) has $\dim K=1$. Let
$\sigma_s>0$ denote the associated smallest singular value at a regular
point.

\begin{lemma}[Quantitative Newton-type error bound]\label{lem:newton}
Let $\Phi:\R^d\to\R^q$ be a quadratic map (each component a polynomial of
degree $\le2$), $\Phi(x^*)=0$, $A:=D\Phi(x^*)$ surjective with smallest
singular value $\sigma>0$, and let $M$ be a Lipschitz constant for $D\Phi$
(finite and global, since $D\Phi$ is affine). Then for every $x$ with
$\|x-x^*\|\le\sigma/(8M)$ and $\|\Phi(x)\|\le\sigma^2/(16M)$,
$$\dist(x,\Phi^{-1}(0))\le\frac2\sigma\|\Phi(x)\|.$$
\end{lemma}
\begin{proof}
Since $\Phi$ is quadratic, Taylor's expansion is exact:
$\Phi(y+v)=\Phi(y)+D\Phi(y)v+B(v,v)$ with $B$ symmetric bilinear and
$\|B(v,v)\|=\tfrac12\|(D\Phi(y+v)-D\Phi(y))v\|\le\tfrac M2\|v\|^2$. Run the
modified Newton iteration $z_0=x$, $z_{i+1}=z_i-A^+\Phi(z_i)$, where $A^+$
is the Moore--Penrose pseudoinverse ($AA^+=I_q$ since $A$ is surjective,
$\|A^+\|=1/\sigma$). Write $b_i=\|\Phi(z_i)\|$, $d_i=-A^+\Phi(z_i)$, so
$\|d_i\|\le b_i/\sigma$. Then
$$\Phi(z_{i+1})=\Phi(z_i)-D\Phi(z_i)A^+\Phi(z_i)+B(d_i,d_i)=(A-D\Phi(z_i))A^+\Phi(z_i)+B(d_i,d_i),$$
so, using $\|A-D\Phi(z_i)\|\le M\|z_i-x^*\|$,
$$b_{i+1}\le\frac{M\|z_i-x^*\|}\sigma\,b_i+\frac M{2\sigma^2}b_i^2.\qquad(*)$$
\textbf{Claim:} every $z_i$ satisfies $\|z_i-x^*\|\le\sigma/(4M)$ and
$b_i\le b_02^{-i}$. By induction: given these, $(*)$ gives
$b_{i+1}\le\tfrac14b_i+\tfrac1{32}b_i<\tfrac12b_i$ (using
$b_i\le b_0\le\sigma^2/(16M)$). Summing the geometric increments,
$\sum_i\|d_i\|\le(b_0/\sigma)\sum_i2^{-i}=2b_0/\sigma\le\sigma/(8M)$, so
$\|z_{i+1}-x^*\|\le\|x-x^*\|+2b_0/\sigma\le\sigma/(4M)$, closing the
induction. The $z_i$ are Cauchy with limit $z_\infty$; $b_i\to0$ and
continuity of $\Phi$ give $\Phi(z_\infty)=0$, and $\|z_\infty-x\|\le2b_0/\sigma$.
No convexity is used anywhere in this argument.
\end{proof}

Fix $x^*\in\CAZ(n,N)$ and let $H=\{c\in\R^n:\sum_kc_k=0\}$, $P_H$ the
orthogonal projection onto $H$. Define the \textbf{balanced constraint map}
$\Phi_s(x)=\bigl(g(x),\,P_Hh(x)/n\bigr)\in\R^{Nn}\times H$, where
$g_{j,m}(x)=|x_j(m)|^2-1$ and $h_k(x)=S_x(k)-Nn$ (so $g(x)=0$ is the
amplitude constraint and $h(x)=0$ is the balanced spectral constraint).

\begin{lemma}[Zero set]\label{lem:phis-zero}
$\Phi_s^{-1}(0)=\CAZ(n,N)$, globally.
\end{lemma}
\begin{proof}
If $g(x)=0$ then $x$ is unimodular, so by Parseval
$\sum_kh_k(x)=n\|x\|^2-nNn=0$, i.e.\ $h(x)\in H$; then $P_Hh(x)=0$ forces
$h(x)=0$, so $x\in\CAZ(n,N)$. Conversely $\CAZ(n,N)\subseteq\{g=0,h=0\}\subseteq\Phi_s^{-1}(0)$.
\end{proof}

\begin{lemma}[Uniform derivative bounds]\label{lem:phis-lip}
$D\Phi_s$ is $2\sqrt2$-Lipschitz globally, and at any $x^*\in\CAZ(n,N)$,
$\|D\Phi_s(x^*)\|_{\mathrm{op}}\le2\sqrt{N+1}$. Both bounds are independent
of $n$.
\end{lemma}
\begin{proof}
\emph{(Lipschitz.)} The rows of $Dg$ are $2x_j(m)$, supported on disjoint
real $2$-planes, so $\|Dg(x)-Dg(y)\|_F=2\|x-y\|$. The $h_k$-row of
$Dh/n$ in block $j$ is (the real form of) $2\hat x_j(k)f_k/n$
($f_k(m)=e^{2\pi ikm/n}$), so by Parseval
$\|Dh(x)/n-Dh(y)/n\|_F^2=4\|x-y\|^2$; $P_H$ is a contraction. Stacking:
Lipschitz constant $\le\sqrt{4+4}=2\sqrt2$.
\emph{(Operator norm at $x^*$.)} $Dg(x^*)Dg(x^*)^\top=4I$ (disjoint unit
$2$-planes, $|x^*_j(m)|=1$); the Gram matrix of the unprojected scaled
spectral rows is $4N\delta_{kk'}$ (Parseval per block, $S_{x^*}\equiv Nn$).
The two blocks have operator norms $2$ and $2\sqrt N$; the stack has norm
$\le\sqrt{4+4N}=2\sqrt{N+1}$.
\end{proof}

\textbf{Hypothesis (R) restated exactly.} $D\Phi_s(x^*)$ is surjective onto
$\R^{Nn}\times H$ iff $\dim K=1$, where $K$ is the dependency space of
Section~\ref{sec:zc}: a pair $(a,c)$ with $c\in H$ annihilates the row
space of $D\Phi_s(x^*)$ iff $\sum_{j,m}a_{j,m}Dg_{j,m}+\sum_k(c_k/n)Dh_k=0$,
and by the definition of $K$ these solutions are parametrized by $c\in K$
(with $a$ determined by $c$). The Parseval solution $c=(1,\dots,1)\notin H$,
so the annihilator inside $\R^{Nn}\times H$ has dimension $\dim(K\cap H)=\dim K-1$
($K$ always contains the constant vector, which is not in $H$); hence
surjectivity $\iff\dim K=1$, matching hypothesis (R) as stated above.
Write $\sigma_s>0$ for the smallest singular value of $D\Phi_s(x^*)$ when
(R) holds, and $\rho:=\sigma_s/(16\sqrt2)$ (so $\sigma_s/(8M)=\rho$ with
$M=2\sqrt2$, Lemma~\ref{lem:phis-lip}).

\begin{theorem}[Theorem C --- Proved conditional on hypothesis (R)]\label{thm:C}
Near a regular $x^*\in\CAZ(n,N)$:
\begin{enumerate}
\item[\rm(a)] \emph{Sharp (linear) growth:}
$D(x)\ge\dfrac{\sigma_s}{2\sqrt2}\dist(x,\CAZ(n,N))$ for $x\in C_1$ near
$x^*$, where $D(x)=Nn-\|x\|^2$.
\item[\rm(b)] \emph{KL exponent $0$:}
$\dist(0,\partial F(x))\ge c_0:=\sigma_s/(8\sqrt2)$ for all non-CAZAC $x$
in a neighborhood of $x^*$.
\item[\rm(c)] \emph{Finite exact arrival:} if the Theorem~A limit $\bar x$
is a regular CAZAC point, then $x^\ell=\bar x$ \textbf{exactly} for all
sufficiently large $\ell$. The number of ``large'' steps before exact
arrival is at most $(Nn-\|x^0\|^2)/c_0^2$.
\end{enumerate}
\end{theorem}
\begin{proof}
(a) The objective deficit is simultaneously the $\ell^1$ slack of both
constraint families:
$D(x)=\sum_{j,m}(1-|x_j(m)|^2)=\tfrac1n\sum_k(Nn-S_x(k))$, all terms
$\ge0$; in particular $\|\Phi_s(x)\|=\sqrt{\|g(x)\|^2+\|P_Hh(x)/n\|^2}\le\sqrt{D(x)^2+D(x)^2}=\sqrt2D(x)$.
\emph{Case 1: $\sqrt2D(x)\le\sigma_s^2/(16M)$, $M=2\sqrt2$.} Apply
Lemma~\ref{lem:newton} to $\Phi=\Phi_s$ (quadratic; zero set $\CAZ(n,N)$ by
Lemma~\ref{lem:phis-zero}; surjective at $x^*$ by (R)):
$\dist(x,\CAZ(n,N))\le(2/\sigma_s)\|\Phi_s(x)\|\le(2\sqrt2/\sigma_s)D(x)$,
i.e.\ $D(x)\ge(\sigma_s/(2\sqrt2))\dist(x,\CAZ(n,N))$.
\emph{Case 2: $\sqrt2D(x)>\sigma_s^2/(16M)$}, i.e.\ $D(x)>\sigma_s^2/64$.
Since $x$ is in the stated neighborhood, $\dist(x,\CAZ(n,N))\le\|x-x^*\|\le\rho=\sigma_s/(16\sqrt2)$,
so $(\sigma_s/(2\sqrt2))\dist(x,\CAZ(n,N))\le\sigma_s^2/64<D(x)$: the bound
holds directly. Either case gives (a).

(b) Let $d:=\dist(x,\CAZ(n,N))>0$ and pick a nearest $z\in\CAZ(n,N)$
(exists, $\CAZ(n,N)$ compact). Any $w\in\partial F(x)$ has the form
$w=v-x$ with $v\in N_{C_1}(x)$ (the sum rule, as in Theorem~A's proof).
Since $z\in C_1$, the normal-cone inequality gives $\langle v,z-x\rangle_\R\le0$,
hence, using $\|z\|^2=Nn$,
$$\langle w,z-x\rangle_\R\le-\langle x,z-x\rangle_\R=\|x\|^2-\langle x,z\rangle_\R
=-\tfrac12(Nn-\|x\|^2)+\tfrac12d^2=-(F(x)-F(x^*))+\tfrac{d^2}2,$$
so Cauchy--Schwarz on the left gives $\|w\|\,d\ge(F(x)-F(x^*))-d^2/2$.
By (a) (in the form $F(x)-F(x^*)\ge(\sigma_s/(4\sqrt2))d$, i.e.\ (a) with
$D(x)=2(F(x)-F(x^*))$), dividing by $d$:
$\|w\|\ge(F(x)-F(x^*))/d-d/2\ge\sigma_s/(4\sqrt2)-d/2$. Since
$d\le\|x-x^*\|\le\rho=\sigma_s/(16\sqrt2)$, this gives
$\|w\|\ge\sigma_s/(4\sqrt2)-\sigma_s/(32\sqrt2)=(7/32)\sigma_s/\sqrt2\ge\sigma_s/(8\sqrt2)=c_0$.
Taking the infimum over $w\in\partial F(x)$ proves (b).

(c) \emph{(Full chain.)} Let $\bar x$ be the Theorem~A limit, regular with
constants $\rho,\sigma_s,c_0$ as in (a)/(b). Since $x^\ell\to\bar x$ and
$\|x^{\ell+1}-x^\ell\|\to0$ (Theorem~A's square-summability), fix $L_0$ so
that for all $\ell\ge L_0$: $x^{\ell+1}$ lies in the (b)-neighborhood of
$\bar x$, and $\|x^{\ell+1}-x^\ell\|<c_0$. By the exact identity from
Theorem~A's proof, $w^{\ell+1}:=x^\ell-x^{\ell+1}\in\partial F(x^{\ell+1})$
with $\|w^{\ell+1}\|=\|x^{\ell+1}-x^\ell\|$ \emph{exactly}, at every $\ell$
-- no approximation. Suppose toward contradiction $x^{\ell+1}\notin\CAZ(n,N)$
for some $\ell\ge L_0$. Then $x^{\ell+1}$ is a non-CAZAC point inside the
(b)-neighborhood, so (b) gives $\dist(0,\partial F(x^{\ell+1}))\ge c_0$. But
$w^{\ell+1}\in\partial F(x^{\ell+1})$ and $\|w^{\ell+1}\|<c_0$ by choice of
$L_0$ --- a contradiction, since $w^{\ell+1}$ itself witnesses a subgradient
of norm $<c_0$. Hence $x^{\ell+1}\in\CAZ(n,N)$ \textbf{exactly}: the
identity is exact and (b)'s bound is a hard floor on every non-CAZAC point
in the neighborhood, so a subgradient below that floor cannot occur unless
the point already lies on $\CAZ(n,N)$. Write $z:=x^{\ell+1}\in\CAZ(n,N)$.
The next step maximizes $\langle z,y\rangle_\R$ over $y\in C_1$; by
Cauchy--Schwarz and Prop.~\ref{prop:2} ($\|z\|=\sqrt{Nn}$ is the global
max-norm), $\langle z,y\rangle_\R\le\|z\|\|y\|\le Nn$ with equality iff
$y=z$ --- the argmax is the singleton $\{z\}$, so $x^{\ell+2}=z$, and by
induction $x^m=z$ for all $m>\ell$; in particular $\bar x=z\in\CAZ(n,N)$.
The step-count bound is Theorem~A's
$\sum_\ell\|x^{\ell+1}-x^\ell\|^2\le Nn-\|x^0\|^2$ telescoped against
$c_0^2$ per large step.
\end{proof}

\begin{remark}
This is a genuine proof by contradiction, not a heuristic squeeze: an
exact identity (equality, from Theorem~A) is jointly satisfiable with a
hard lower bound (inequality, from (b)) on the same quantity only at
$\dist=0$. Parts (a), (b), (c) share one status tag because they rest on
exactly the same two inputs --- hypothesis (R) at $\bar x$ (giving
$\sigma_s,\rho$) and the exact H2 identity from Theorem~A's proof --- and
(c) adds no assumption beyond the theorem's own hypothesis
``$\bar x\in\CAZ(n,N)$'' plus the unconditional singleton-argmax fact.
\end{remark}

\begin{remark}[The $N=1$ vs.\ $N=2$ gap is local-rate-invariant]
Theorem C's constants do not distinguish $N=1$ from $N=2$: first-order
sharpness at a regular point decides the local geometry identically. The
empirically observed difference in success probability between the
single-sequence and coupled formulations is therefore a global
(basin-of-attraction) phenomenon, not a local-rate one.
\end{remark}

\begin{remark}[Novelty check against generic identifiability theory]
Neither the Lewis--Luke--Malick theory of local linear convergence for
averaged/alternating projections \cite{LLM09} nor the Bonnans--Shapiro /
Drusvyatskiy--Lewis generic identifiability program \cite{DL12} delivers
Theorem~C. The former's core hypothesis, strongly regular (transversal)
intersection, is exactly what Section~\ref{sec:zc} shows \emph{fails} at
every CAZAC point (the universal Parseval dependency), and even where
satisfiable its conclusion is R-linear convergence of a different
algorithm, never finite exact arrival. The latter's genericity condition is
plausibly satisfied at CAZAC points, but its machinery yields only a
qualitative locally-minimal-identifiable-set statement, not a quantitative
sharp-growth constant, KL exponent, or finite-step termination guarantee.
\end{remark}

\section{Theorem D and regularity (R)}\label{sec:zc}

Write $x_u(m)=\omega^{-uT_m}$ ($\omega=e^{2\pi i/n}$, $T_m=m(m+1)/2$, $n$
odd) or $x_u(m)=\omega_{2n}^{-um^2}$ ($\omega_{2n}=e^{i\pi/n}$, $n$ even)
for the Zadoff--Chu sequence of root $u$, $\gcd(u,n)=1$ \cite{Chu72}. At a
CAZAC tuple all constraints are active; write
$$K=\{c\in\R^n: \overline{x_j(m)}\,\bigl(F^*(c\odot\widehat x_j)\bigr)(m)\in\R\ \ \forall j,m\}$$
for the active-gradient dependency space (the tangent space of the
constraint intersection has dimension $(N-1)n+\dim K$; hypothesis (R) is
$\dim K=1$).

\begin{theorem}[Theorem D --- Proved for all $n\ge1$, odd and even]\label{thm:D}
Let $n=\prod_p p^{a_p}$, $d_{\max}=\prod_p p^{\lfloor a_p/2\rfloor}$ (the
largest $d$ with $d^2\mid n$). Then for every $n\ge1$, every $N\ge1$, every
tuple of coprime roots $u_1,\dots,u_N$,
$$K=\{c\in\R^n: c(k+d_{\max})=c(k)\ \forall k\},\qquad \dim_\R K=d_{\max},$$
independent of $N$, of the roots, and of the parity of $n$.
\end{theorem}
\begin{proof}
\textbf{Step 0.} A dependency $\sum_{j,m}a_{j,m}\nabla g_{j,m}+\sum_kc_k\nabla h_k=0$
evaluated on a perturbation is equivalent to $F^*(c\odot\widehat x_j)$
being coordinatewise a real multiple of $x_j$ --- the defining condition
of $K$; dependencies with $c=0$ force $a=0$ since amplitude gradients have
disjoint supports. So the dependency space is (isomorphic to) $K$.

\textbf{Step 1 (self-transform).} For odd $n$, writing $k=us$,
$\widehat x_u(k)=G_u\,\omega^{uT_{u^{-1}k}}$ with $G_u=\sum_j\omega^{-uT_j}$,
$|G_u|^2=n$ by Parseval, using the exact identity $T_{m+s}=T_m+T_s+ms$. For
even $n$, the parallel identity holds with $T_m$ replaced by $m^2$,
exponents tracked mod $2n$: $\widehat x_u(k)=G_u\,\omega_{2n}^{us^2}$,
using $\omega_{2n}^{v(j+n)^2}=\omega_{2n}^{vj^2}$ (an exact integer
identity valid because $n$ is even).

\textbf{Step 2 (reduction to a correlation equation).} Substituting Step 1
into the dependency condition and reindexing $k=us\bmod n$ turns
$c\in K$ into reality, for all $m$, of a convolution
$\sum_sc'_s\,e(s+m)$, where $e$ is built from $G_u$ and the same phase.

\textbf{Step 3 (Fourier solve).} Taking the length-$n$ DFT in $m$ reduces
this to $\hat c'(t)=0$ wherever $\hat e(-t)\ne0$. In both parities,
$\hat e(t)=0$ exactly when $\sigma^2\equiv0\pmod n$ ($\sigma=u^{-1}t$): the
zero set is the subgroup $Z_0=\{t: t^2\equiv0\bmod n\}=m_0\Z_n$,
$|Z_0|=n/m_0=d_{\max}$ where $m_0=\prod_pp^{\lceil a_p/2\rceil}$.

\textbf{Step 4.} $\hat c'(t)=\hat c(u^{-1}t)$, so ``$\operatorname{supp}\hat c'\subseteq Z_0$''
is the same condition for every root and block, giving
$K=\{c:\operatorname{supp}\hat c\subseteq m_0\Z_n\}$, i.e.\ the
$d_{\max}$-periodic real vectors, of real dimension $d_{\max}$. This count
is read off directly from $|Z_0|=d_{\max}$ and is insensitive to parity:
an alternate derivation via counting conjugate-symmetric pairs
$\{t,-t\bmod n\}$ in $Z_0$ would need to treat the self-paired frequency
$t=n/2$ separately when $n\equiv0\bmod4$ places $n/2\in Z_0$, but that
bookkeeping is specific to the pair-counting route and does not arise
here, since Steps 1--4 characterize $K$ by its Fourier support directly.
\end{proof}

\begin{corollary}\label{cor:D1}
Hypothesis (R) holds at every Zadoff--Chu tuple iff $n$ is squarefree ---
in particular at every prime $n$, odd or $n=2$. Theorem~C is therefore
\textbf{unconditional} whenever the Theorem~A limit lies on the symmetry
orbit of a Zadoff--Chu tuple and $n$ is squarefree.
\end{corollary}

\begin{corollary}\label{cor:D2}
At Zadoff--Chu tuples with squarefree $n$ (either parity), the CAZAC
tangent dimension is exactly $(N-1)n+1$: a $1$-dimensional phase circle for
$N=1$ versus an $(n+1)$-dimensional manifold for $N=2$.
\end{corollary}

Numerical verification: {\tt experiments/verify\_zc\_regularity.py} for odd
$n\in\{3,5,7,9,15,21,25,27,45,49,75\}$ (21 configurations) and
{\tt experiments/verify\_even\_n\_regularity.py} for the $13$ even values
$n\in\{2,4,6,8,10,12,14,16,18,20,24,36,50\}$ (43 configurations); all match
$\dim K=d_{\max}$ exactly. No Chebotar\"ev-type nonvanishing-minors
argument is needed: the DFT of a chirp is a chirp, collapsing the
dependency analysis to counting solutions of $t^2\equiv0\bmod n$. Both
scripts are in the companion repository \cite{CazacAlgorithmRepo}.

\subsection{The chirp-argument schema, and where it does and does not generalize}

Stripped of Zadoff--Chu-specific notation, Theorem~D's proof needs exactly:
(1) a closure property, $\widehat{x_u}=G_u\cdot x_u\circ\phi_u$ for a
scalar $G_u$ and an index-set automorphism $\phi_u$; (2) an algebraic
addition law for the phase making the key substitution exact; (3) applying
the closure twice, once to reduce the dependency condition to a
correlation equation, once more to diagonalize that equation into a
congruence on the Fourier support of the kernel. This replaces a
linear-algebra rank question with a congruence-counting question, and is a
genuinely reusable schema, not a one-off computation.

\begin{proposition}[Positive lead --- Proved classically, newly connected here]
The Legendre symbol $\chi$ modulo an odd prime $p$ satisfies
$\hat\chi(k)=\chi(k)\,g(\chi)$ for $k\ne0$ ($g(\chi)$ the Gauss sum,
$|g(\chi)|^2=p$): $\chi$ is, up to a scalar, its own DFT, the same closure
shape used by Theorem~D, with trivial reindexing.
\end{proposition}

\begin{proposition}[Negative result --- Proved, standard exponential-sum
theory, newly connected here]
Cubic and higher-degree phase (``algebraic chirp'') sequences do not have
this closure property: quadratic Gauss sums have an exact closed-form
evaluation (magnitude $\sqrt n$ always), which is what underwrites
Theorem~D's Step~1, but Weyl sums of degree $\ge3$ have no such exact
evaluation, only upper bounds, and the addition law
$(m+s)^d=m^d+s^d+(\text{nonlinear mixing})$ for $d\ge3$ does not reduce to
a reindexing of the original phase. \textbf{No variant of the chirp
argument applies to cubic or higher phase sequences.}
\end{proposition}

\subsection{The Bj\"orck case: a partial result, not a proof}

Björck's construction (odd prime $p$, Legendre symbol $\chi$)
\cite{Bjorck90} defines $b(m)=e^{i\theta(m)}$ with $\theta$ a two-valued
function of $\chi(m)$, differing by $p\bmod4$. Unlike $\chi$ itself, $b$
is not literally self-dual under the DFT, but a different closure fact
rescues part of the argument.

\begin{lemma}[Proved, exact symbolic computation]\label{lem:bjaffine}
$b$ is exactly \textbf{affine} in $\chi$: $b=A\cdot\mathbf1+B\cdot\chi+C\cdot\delta_0$
with $C=1-A$, for explicit constants $A,B$ depending on $p\bmod4$. The
$3$-dimensional space $T:=\operatorname{span}_\C\{\mathbf1,\chi,\delta_0\}$
closes under the DFT: $\hat b=C\cdot\mathbf1+Bg\cdot\chi+Ap\cdot\delta_0$.
\end{lemma}

\begin{lemma}[Proved]\label{lem:bjequiv}
Let $G$ be the group of quadratic-residue units mod $p$ (order $q=(p-1)/2$),
acting by decimation. $b,\hat b$ are pointwise $G$-invariant, so $K$ is a
$G$-invariant subspace of $\R^p$. $\R^p\cong\mathbf1_{\{0\}}\oplus\R[G]\oplus\R[G]$
as a $G$-module, with trivial-isotypic piece exactly $T\cap\R^p$
(dimension $3$).
\end{lemma}

\begin{theorem}[Proved, both residue classes]\label{thm:E0}
$K\cap T=\operatorname{span}\{\mathbf1\}$; i.e.\ restricted to the
$3$-dimensional symmetry-invariant candidate space, the only dependency is
the Parseval one, for every prime $p$.
\end{theorem}
\begin{proof}
Restricting the dependency condition
$\overline{b(m)}\,(F^*(c\odot\hat b))(m)\in\R$ for all $m$ to
$c=\alpha\mathbf1+\beta\chi+\gamma\delta_0\in T$ collapses --- since
$b,\hat b\in T$ closes under pointwise product and $F$ --- to exactly
three scalar equations, one evaluation each at $m=0$, $m\in$~QR,
$m\in$~QNR (the condition depends on $m$ only through $\chi(m)\in\{0,1,-1\}$).
Write $\alpha_0$ (resp.\ $\beta_0$) for the phase such that $Bg=s\sqrt p\,e^{i\alpha_0}$
with $s=\pm1$ when $p\equiv1\pmod4$ (resp.\ the corresponding phase when
$p\equiv3\pmod4$), using that the Gauss sum $g=g(\chi)$ is real for
$p\equiv1\pmod4$ and purely imaginary ($g=i\sqrt p$, up to sign) for
$p\equiv3\pmod4$ --- this parity of $g$ is exactly where the two residue
classes diverge. The three equations reduce, after eliminating the
$m=0$ equation (which involves only $\alpha,\gamma$ and is solved by
$\alpha$ free), to:
$$p\equiv1\!\!\pmod4:\quad s\sqrt p\,\beta=0,\qquad -s\gamma\bigl(1+c(p-1)\bigr)=0\quad(c=\cos\alpha_0,\,s=\sin\alpha_0>0);$$
$$p\equiv3\!\!\pmod4:\quad \frac{\sqrt p(1-\cos\beta_0)}2\,\beta=\frac{p\sin\beta_0}2\,\gamma,
\qquad \frac{\sqrt p(1-\cos\beta_0)}2\,\beta=-\frac{\sin\beta_0}2\,\gamma.$$
In the first case, $s\sqrt p\ne0$ forces $\beta=0$ directly, and
$1+c(p-1)>0$ (as $c>-1$ and $p\ge5$) then forces $\gamma=0$. In the second
case, subtracting the two equations gives
$\gamma\bigl(\tfrac{p\sin\beta_0}2+\tfrac{\sin\beta_0}2\bigr)=0$, i.e.\
$\gamma\sin\beta_0(p+1)/2=0$; since $p+1\ne0$ and $\sin\beta_0\ne0$ (a
degenerate $\beta_0\in\{0,\pi\}$ would force $g$ real, contradicting
$g=i\sqrt p$), $\gamma=0$, and then $\beta=0$ from either equation
(the coefficient $\sqrt p(1-\cos\beta_0)/2$ is nonzero unless
$\beta_0=0$, again excluded). In both residue classes the system forces
$\beta=\gamma=0$ with $\alpha$ free, giving $K\cap T=\operatorname{span}\{\mathbf1\}$
exactly.
\end{proof}

\begin{remark}[What is proved and what is not --- do not upgrade]
Lemma~\ref{lem:bjaffine} and Lemma~\ref{lem:bjequiv} and
Theorem~\ref{thm:E0} are complete proofs. What remains \textbf{unproved} is
whether $K$ meets any of the $q-1$ nontrivial isotypic pieces of $\R^p$
under $G$: by Schur's lemma the dependency operator restricted to each
such (complex $2$-dimensional) block is a $2\times2$ matrix, structurally a
Jacobi-sum-type bilinear expression, and its nonvanishing was not proved.
A second, independent attempt at a clean closed-form nonvanishing argument
also failed: the relevant determinants are confirmed numerically nonzero
at every one of $44$ tested primes $p\in[5,200)$ (both residue classes,
$\dim K=1$ throughout), but with no uniform lower bound of the shape a
classical single-Gauss/Jacobi-sum-modulus argument would give --- at
$p=71$ the relevant gap is $\sim10^4\times$ smaller than at neighboring
primes, a genuine near-degeneracy confirmed not to be a rounding artifact
by an independent recomputation at $60$-digit precision at that specific
prime. Script: {\tt experiments/verify\_bjorck\_isotypic.py} in the
companion repository \cite{CazacAlgorithmRepo} (the $44$-prime
determinant check; a separate, disjoint set of $14$ primes is used for
the affine-closure check of Lemma~\ref{lem:bjaffine}, in
{\tt experiments/verify\_bjorck\_regularity.py}). \textbf{The overall regularity
claim at Bj\"orck points is therefore still open, not proved}, despite this
genuine partial progress: what is established is a full proof restricted
to the symmetry-invariant subspace, plus strong (but non-uniform)
numerical evidence, not a proof, for the complementary directions.
\end{remark}

\section{Open problems}

\noindent\textbf{Q1 (settled).} Does $x^\ell$ converge to a single point? Yes ---
Theorem~\ref{thm:A}.

\noindent\textbf{Q2 (open, reframed).} Characterize all fixed points of $(\star)$
for general $n,N$. The stationarity system is a KKT eigenvalue equation
$x^*=(\Lambda+F^*MF)x^*$ for diagonal multipliers, structurally reminiscent
of discrete prolate/Slepian time-and-band-limiting operators. Complete for
$n=2$ (four symmetry classes); open in general. A sharper sub-question:
are non-CAZAC fixed points always measure-zero repellers, or can they
attract an open set of initializations?

\noindent\textbf{Q3 (partially answered).} Why does $N=2$ empirically always
succeed where $N=1$ sometimes fails? The dimension lower bound
$\dim\CAZ(n,N)\ge(N-1)n+1$ (proved) gives a mechanism, and
Corollary~\ref{cor:D2} makes it exact at squarefree Zadoff--Chu points; but
the pathological fixed-point classes are the same for $N=1,2$, so the gap
is about basin measure, not existence, and remains open.

\noindent\textbf{Q4 (essentially settled locally; one remaining gap).} Rate of
convergence near $\CAZ(n,N)$: settled by Theorem~C at regular points, and
Theorem~D now makes regularity unconditional at Zadoff--Chu points of
either parity for squarefree $n$. Remaining: (R) at non-Zadoff--Chu CAZAC
points is open beyond the Bj\"orck partial result above; $n$-uniformity of
the true sharp-growth constant is conjectural (proved bounds scale like
$c_N/n$, sampled ratios show no decay).

\noindent\textbf{Q5 (open, deprioritized).} What does the $\tau$-continuation
presolve actually buy? Hypothesis: it drives the iterate onto (or near) the
easier flat-spectrum variety before the main phase. Untested by ablation.

\noindent\textbf{Q6 (the central open question).} For $N\ge2$ and Gaussian
initialization, does the iteration converge to $\CAZ(n,N)$ with
probability $1$ (or $\ge1-c^n$)? Requires a full fixed-point census (Q2)
and a basin/measure argument. For $N=1$, empirical success probability
decays from $\approx1$ at small $n$ to $\approx0.3$ at $n=167$; whether
this decay is exponential or polynomial in $n$ is unresolved.

\noindent\textbf{Q7 (open, newly raised).} Does every trajectory of the elliptope
iteration converge to a rank-1 point? Theorem~\ref{thm:ell} proves
single-point convergence unconditionally but not convergence to an extreme
point. Conjecture: for a.e.\ (e.g.\ Gaussian-perturbed) initialization on
$\El$, attracting fixed points are exactly the rank-1 extreme points, and
the higher-rank fixed-point continuum (containing $I_n$) is a measure-zero
repelling set. This is structurally the same conjecture as Q6 in a second,
unrelated domain: if true, it is independent evidence for Q6's general
shape; if false, CAZAC's extra structure (the dimension count of
Corollary~\ref{cor:D2}) may be doing more work than currently credited.
\emph{Attack plan:} extend Prop.~\ref{prop:ell4}'s experiment to
$n\in\{5,10,20\}$ with many Gaussian-perturbed starting points per $n$, and
histogram the rank of the limit point; a rank-1 outcome in every trial
would support attempting a basin/measure argument analogous to Q6's, while
a robust fraction of non-rank-1 limits would falsify the conjecture as
stated and make the counterexamples themselves the next object of study.

\section{Verification status and provenance}

\textbf{Tiers.} \emph{Unconditional}: complete proof, no open hypothesis.
\emph{Proved conditional on (R)}: complete proof, resting on regularity
hypothesis (R). \emph{Proved conditional on (R) and an open sub-case}:
complete proof modulo one explicitly identified, unresolved lemma.
\emph{Verified numerically}: no proof; a script reproduces the claim to a
stated tolerance. \emph{Conjecture}: stated as a target, not established
either way.

\textbf{Principal claims.}
\begin{itemize}
\item Propositions~\ref{prop:1}--\ref{prop:4} (generalized power method,
max-norm = CAZAC, monotone ascent, limit points fixed): \emph{Unconditional}.
\item Theorem~A (global single-point convergence): \emph{Unconditional}.
\item Elliptope semialgebraicity and H1--H3 transfer
(Props.~\ref{prop:ell1}--\ref{prop:ell2}): \emph{Unconditional}. Fixed-point
continuum (Prop.~\ref{prop:ell3}): \emph{Unconditional} (cited fact plus a
machine-exact identity). Single-point convergence, generic escape to
rank-1 (Prop.~\ref{prop:ell4}): \emph{Verified numerically} ($n=5$, four
trajectories). Universal rank-1 convergence (Q7): \emph{Conjecture}.
\item Theorem B (obstruction, exposed non-CAZAC fixed point):
\emph{Unconditional}.
\item Theorem C, parts (a),(b),(c), including the full contradiction-based
proof of part (c): \emph{Proved conditional on hypothesis (R)}; all three
parts carry the same status tag.
\item Theorem D (regularity of Zadoff--Chu points, $\dim K=d_{\max}$),
\textbf{both odd and even $n$}: \emph{Unconditional}, a complete proof for
all $n$ regardless of parity.
Corollaries~\ref{cor:D1}--\ref{cor:D2}: \emph{Unconditional} (Corollary~1
conditionally activates Theorem~C, but is itself unconditional as a
dimension/regularity statement).
\item Legendre-symbol closure positive lead: classical fact, newly
connected to the schema, \emph{Unconditional} as stated (no claim of
Bj\"orck regularity is made from it alone). Cubic/higher-phase negative
result: \emph{Unconditional} (a genuine proved negative, not a
placeholder).
\item Bj\"orck partial result: Lemma~\ref{lem:bjaffine},
Lemma~\ref{lem:bjequiv}, Theorem~\ref{thm:E0} ($K\cap T=\operatorname{span}\{\mathbf1\}$):
\emph{Unconditional}. Full regularity (R) at Bj\"orck points:
\textbf{open}, reduced to one unproved determinant-nonvanishing lemma,
\emph{verified numerically only} at 44 primes, $p\in[5,200)$. Do not read
this as ``proved'' --- the reduction is real but incomplete.
\item Q1--Q7: precise open-problem statements as in Section~6; Q1 is
settled (restated as a theorem above), the rest range from ``partially
answered'' to fully open conjectures, as stated there.
\end{itemize}

\textbf{Provenance.} This document was written with substantial assistance
from large language models (Claude, Anthropic) under the direction of
Jeffery Kline, drawing on a longer working record of derivations kept in
the companion repository \cite{CazacAlgorithmRepo}'s {\tt theory/}
directory; every numerical claim above is backed by a named script in
that repository's {\tt experiments/} directory. Every claim above carries
the status tag most consistent with that record; none has been upgraded
in the writing of this paper.

\begin{thebibliography}{99}

\bibitem{ABS13} H. Attouch, J. Bolte, B. F. Svaiter, \emph{Convergence of
descent methods for semi-algebraic and tame problems: proximal algorithms,
forward--backward splitting, and regularized Gauss--Seidel methods},
Math. Program. \textbf{137} (2013), 91--129.

\bibitem{BDLS07} J. Bolte, A. Daniilidis, A. S. Lewis, M. Shiota,
\emph{Clarke subgradients of stratifiable functions}, SIAM J. Optim.
\textbf{18}(2) (2007), 556--572. arXiv:math/0601530.

\bibitem{FKP21} P. Felzenszwalb, C. Klivans, A. Paul, \emph{Iterated Linear
Optimization}, Quarterly of Applied Mathematics (2021). arXiv:2012.02213.

\bibitem{JNRS10} M. Journ\'ee, Y. Nesterov, P. Richt\'arik, R. Sepulchre,
\emph{Generalized Power Method for Sparse PCA}, J. Mach. Learn. Res.
\textbf{11} (2010), 517--553. arXiv:0811.4724.

\bibitem{LLM09} A. S. Lewis, D. R. Luke, J. Malick, \emph{Local linear
convergence for alternating and averaged nonconvex projections}, Found.
Comput. Math. \textbf{9}(4) (2009), 485--513. arXiv:0709.0109.

\bibitem{DL12} D. Drusvyatskiy, A. S. Lewis, \emph{Optimality,
identifiability, and sensitivity}, arXiv:1207.6628.

\bibitem{Chu72} D. C. Chu, \emph{Polyphase codes with good periodic
correlation properties}, IEEE Trans. Inform. Theory \textbf{18}(4) (1972),
531--532.

\bibitem{Bjorck90} G. Bj\"orck, \emph{Functions of modulus $1$ on $\Z_n$
whose Fourier transforms have constant modulus, and ``cyclic $n$-roots''},
in: Recent Advances in Fourier Analysis and its Applications, NATO ASI
Series C \textbf{315}, Kluwer, 1990, 131--140.

\bibitem{CazacAlgorithmRepo} J. Kline, \emph{cazac-algorithm}: companion
research repository holding the {\tt theory/} derivations, the dated
{\tt AUDIT-LEDGER.md}, and the {\tt experiments/} scripts (including
{\tt verify\_zc\_regularity.py}, {\tt verify\_even\_n\_regularity.py},
{\tt verify\_elliptope\_convergence.py}, {\tt verify\_bjorck\_regularity.py},
{\tt verify\_bjorck\_isotypic.py}) that reproduce every numerical claim in
this paper, 2026.

\end{thebibliography}

\end{document}
